\chapter{Conclusion}\label{conclusion}

This diploma project set out to design, implement and evaluate a complete software system capable of automatically detecting individual trees in satellite imagery of the city of Astana. The work was motivated by the practical need of \emph{Zelenstroy} and other municipal services for a fast, repeatable and inexpensive alternative to the manual field survey that has historically been the only way to obtain an urban-tree inventory at the scale of a major city. The analysis of the existing literature (Chapter 1) confirmed that the deep-learning toolkit needed for such a system is technically mature --- state-of-the-art object-detection and instance-segmentation networks routinely exceed 90 \% mean average precision on dedicated public datasets --- but also that \textbf{no published work} to date has evaluated any of these methods on Central-Asian satellite imagery in general or on Astana imagery in particular. Closing this gap was the principal motivation of the present work.

\section{Achievement of the stated objectives}\label{achievement-of-the-stated-objectives}

The six objectives formulated in the introduction were addressed as follows.

\textbf{Objective 1 --- survey the deep-learning state of the art.} Chapter 1 contains a critical analysis of 31 peer-reviewed publications of 2019 -- 2025, organised by technical paradigm (two-stage detectors including Mask R-CNN, one-stage detectors of the YOLO family, semantic segmentation, the domain-specialised DeepForest model and the SAM / SAM 2 foundation models) and culminating in a consolidated Table 1.2 of the best metrics reported by each work. The survey established two essential facts: that fine-tuned models routinely achieve F1-scores in the 0.65 -- 0.85 range on European and American urban data, and that the geographic-generalisation gap between an off-the-shelf and a fine-tuned model is in the order of 0.30 F1 points.

\textbf{Objective 2 --- quantify the gap on Astana imagery.} Section 3.4.1 reports the off-the-shelf DeepForest baseline on Astana --- precision ≈ 0.72, recall ≈ 0.58 --- confirming both the order-of-magnitude consistency with the European urban literature and the necessity of a fine-tuning stage.

\textbf{Objective 3 --- build a custom annotated Astana dataset.} Section 3.2 documents two iterations of dataset construction: a version-1 dataset of 20 source images / 2 242 polygons obtained through from-scratch manual annotation in CVAT, and a version-2 dataset of 77 source images / ≈ 8 000 polygons obtained through a model-in-the-loop pre-labelling workflow that reduced the annotation cost per image by approximately 70 \%. The dataset is, to the best of the authors' knowledge, the first of its kind for Astana and is small but reusable for future research.

\textbf{Objective 4 --- train and combine four complementary models.} Chapters 2 and 3 document the training of three successive YOLOv8x-seg checkpoints on the Astana dataset: the version-1 model (397 epochs on 20 source images, Box mAP@50 = 0.478 on its own small 4-tile validation set), the version-2-from-scratch model (204 epochs from COCO weights on the merged 77-image dataset) and the version-2-finetune model (continued training of the v1 checkpoint on the 57 new images only, 99 epochs). All three models were re-evaluated on a common, larger and harder version-2 validation set; on this like-for-like comparison the version-2-finetune model is the best, with Box mAP@50 = \textbf{0.372} and Mask mAP@50 = \textbf{0.331}, representing a 40 \% relative improvement over the version-1 baseline (0.265 / 0.240). The DeepForest RetinaNet detector, owned by team member Anuar Totin and trained on an independent bounding-box-level annotation set, reaches precision ≈ 0.80 and recall ≈ 0.65 on the Astana validation tiles. A Mask R-CNN instance segmenter, owned by team member Berik Sharipov, is trained on the same polygon dataset as the YOLO branch and provides a direct one-stage vs two-stage architectural comparison (numerical results in Section 3.4). SAM 2 is integrated as a zero-shot mask-refinement stage on top of DeepForest. The YOLO and DeepForest branches are combined through a Weighted-Box-Fusion ensemble (Box mAP@50 ≈ 0.51 on the v1 val), exploiting the complementary failure modes of the two architectures: YOLO over-segments dense canopies, DeepForest under-segments them.

A methodological observation of the version-2 training comparison is worth retaining for future work. A first ablation, restricted to the comparison of v1 against v2-from-scratch, suggested that fine-tuning offered no benefit; this conclusion was reversed once a third configuration --- fine-tune on the new-images subset rather than on the full merged corpus --- was added to the comparison. The episode is a useful reminder that conclusions about transfer-learning effectiveness are extremely sensitive to the precise composition of the fine-tuning set, and that a two-way ablation can hide an important interaction with a third design dimension.

\textbf{Objective 5 --- design and implement an end-to-end software pipeline.} Chapter 2 describes the resulting system: a FastAPI Python backend with a pluggable model-adapter interface, a React 18 single-page frontend with a Leaflet map, a four-mode geographic-conversion module supporting GeoTIFFs and informal screenshots alike, an in-browser map-capture feature that stitches ESRI World Imagery tiles for an arbitrary user-drawn rectangle, and three export formats (GeoJSON, CSV, standalone HTML) that integrate the inventory with the user's existing GIS workflow. The complete prototype is approximately 6 000 lines of Python and JavaScript.

\textbf{Objective 6 --- evaluate the system on realistic Astana scenes.} Section 3.7 reports the end-to-end behaviour of the integrated pipeline: a 1 km × 1 km area captured at zoom 18 is processed in approximately 18 seconds on a single laptop GPU, comfortably below the 30-second budget set by the requirements. Qualitative inspection of the predictions on Astana validation tiles confirms that the system already produces an inventory that is qualitatively informative for the \emph{Zelenstroy} end user, while the aggregate metrics --- Box mAP@50 ≈ 0.48 for YOLO alone and ≈ 0.51 for the ensemble --- establish the empirical baseline against which any future model can be compared.

\section{Scientific contributions}\label{scientific-contributions}

The work contributes the following original results.

\begin{enumerate}
\def\labelenumi{\arabic{enumi}.}
\item
  The \textbf{first published evaluation} of state-of-the-art deep-learning tree-detection models on satellite imagery of Astana and, by extension, of any major Central-Asian capital. The empirical numbers --- off-the-shelf DeepForest precision/recall on Astana, fine-tuned DeepForest on Astana, YOLOv8x-seg on Astana, ensemble --- are now part of the empirical record for the region.
\item
  A \textbf{small but reusable annotated dataset} of Astana satellite imagery, with approximately 77 source images and approximately 8 000 polygon-level tree annotations after tiling, distributed under the same convention as the original DeepForest dataset.
\item
  A \textbf{four-model architecture} combining one-stage instance segmentation (YOLOv8-seg), two-stage instance segmentation (Mask R-CNN), bounding-box detection (DeepForest fine-tuned) and zero-shot mask refinement (SAM 2 on DeepForest boxes), with a Weighted-Box-Fusion ensemble of the detector outputs. The complementarity of the YOLO and DeepForest failure modes --- established quantitatively in Section 3.5.3 --- is a novel observation in the urban-tree literature.
\item
  A \textbf{complete deployable prototype} --- backend, frontend, geographic conversion, three exporters, an in-browser ESRI tile-capture feature --- that serves as a template for any city administration wishing to adopt deep-learning-based tree inventories. The system is designed to be retrained on imagery from any other city by simply replacing the dataset.
\end{enumerate}

\section{Practical contributions}\label{practical-contributions}

The system is in a state in which it can be used immediately for an informal Astana tree inventory: the user supplies a satellite image (uploaded or captured interactively from the map), runs the selected model, inspects the result on a Leaflet map and exports the inventory in the format required by the downstream tool. The interactive response time of approximately 18 seconds for a 1 km × 1 km area is short enough that a single specialist can comfortably process several districts of the city in a working day --- a productivity improvement of two to three orders of magnitude over the manual field survey it replaces.

The architectural decisions made along the way --- adapter pattern for the model layer, four-mode geographic conversion, sliding-window tiled inference with global NMS, model-in-the-loop pre-labelling for annotation expansion, the Cyrillic-aware data-preparation pipeline --- are independent of the specific application to Astana and can be reused as a template for any future urban-remote-sensing project.

\section{Limitations and future work}\label{limitations-and-future-work}

The current implementation has the limitations listed in Section 3.9: a validation set of only 4 (v1) or 10 (v2) tiles, a single class label that does not distinguish species, no explicit shadow modelling, an iterative DeepForest fine-tune that should be refactored into a single multi-epoch run, no quantitative SAM 2 / Mask R-CNN / YOLO mask ablation, an in-memory job store with no persistence and a single-laptop deployment.

The natural directions for future work are:

\begin{enumerate}
\def\labelenumi{\arabic{enumi}.}
\item
  \textbf{Dataset expansion to 200 -- 500 source images}, with the explicit goal of reaching the F1 = 0.7 range demonstrated by {[}3{]} and {[}6{]} on European urban data. With more data, all three trained branches (YOLOv8-seg, Mask R-CNN, fine-tuned DeepForest) should converge towards the European literature-baseline performance, and the ensemble of the three should narrow the remaining gap.
\item
  \textbf{Quantitative ablation of the four branches under a single protocol.} The present work reports per-branch numbers but does not yet provide a unified mask-quality ablation that compares YOLO masks, Mask R-CNN masks and SAM 2-refined DeepForest masks on the same set of ground-truth crowns. Such an ablation is required before any production use of either the ensemble or the SAM 2 output.
\item
  \textbf{Species-level annotation and a multi-class detection head}, following the multi-task approach of {[}10{]} and the species-specific work of {[}18{]}. \emph{Zelenstroy} explicitly requested this extension during early-stage discussions, since species-specific pruning and replacement is a separate budget line in the municipal urban-forestry plan.
\item
  \textbf{SAM 2 video propagation for multi-temporal monitoring.} The present work uses only SAM 2's image-level segmentation, but its native streaming-memory mode could be leveraged to track individual crowns across multiple satellite acquisitions of the same district, automatically detecting tree mortality and new plantings between successive captures.
\item
  \textbf{Sentinel-2 integration} for city-wide canopy-cover monitoring, following the methodology of {[}14{]} and the canopy-height extension of {[}16{]}. A Sentinel-2-based monitoring pipeline would complement the very-high-resolution per-tree inventory built in the present work and would allow the city to observe long-term trends in green coverage.
\item
  \textbf{A persistent inventory database} that accumulates detection results over time, supports versioning of the canonical inventory and exposes a difference view (trees added, trees removed) between successive captures of the same district.
\item
  \textbf{Multi-user deployment} through a containerised version of the system, with authentication and per-organisation isolation, so that \emph{Zelenstroy} and other municipal services can adopt the tool as part of their day-to-day operations.
\end{enumerate}

\section{Closing remark}\label{closing-remark}

This work has demonstrated that the deep-learning techniques developed for tree detection in American forests, European cities and Asian metropolises can be adapted to the very different geographic, architectural and floristic context of a Central-Asian capital through a relatively modest annotation effort and an off-the-shelf software stack. The resulting system delivers, for the first time, a quantitative baseline for automated urban-tree detection in Astana, an open prototype that can be operated by a non-technical user from a web browser, and a reusable template for future work in the region. The combination of these three contributions defines the engineering and scientific value of the present diploma project.

\newpage
